\chapter{Organizing a Mind}

\chapterepigraph{%
  The mind is not a library.\\
  It is a living city:\\
  some districts thriving, some decaying,\\
  connected by roads of varying reliability,\\
  always under construction.
}{Flyxion, \textit{Semantic Infrastructure}}

\sidenote{Connects to\\Ch.\,23 memory\\chains and\\Ch.\,32 beyond\\files and folders}

Having established what learning is (infrastructure construction),
what its dynamic is (exploration-exploitation), and what tools extend
it (cognitive prosthetics), this chapter turns to the practical question:
how should a mind organize itself?

\section{Notes}

External notes are memory prosthetics (Chapter~35): they extend the
biological memory chain into physical or digital space. But different
note-taking systems implement different cognitive architectures.

\begin{itemize}
  \item \textbf{Linear notes} (sequential writing): implement a
        1D memory chain. Easy to produce, hard to navigate. Good for
        capturing temporal sequence (lecture notes, journal entries).
  \item \textbf{Hierarchical notes} (outlines, nested folders):
        implement a tree. Better than linear for structure, but imposes
        a single categorization. The same problem as file systems:
        concepts belong to multiple categories.
  \item \textbf{Networked notes} (Zettelkasten, wiki-style):
        implement a graph. Each note is a node; links are edges.
        Supports the full topology of knowledge. Harder to maintain
        but much richer in structure.
  \item \textbf{Spatial notes} (mind maps, concept maps):
        implement a geometric arrangement. Visual proximity encodes
        semantic proximity. Supports intuitive navigation but
        poor at precise retrieval.
\end{itemize}

The networked note system implements the knowledge landscape of
Chapter~34: a graph of concepts linked by semantic relationships,
navigable by following high-accessibility paths.

\section{Memories}

Memories, unlike notes, are not explicitly organized — they are
implicitly organized by the constraint structure they form part of.
But the agent can influence this organization through \emph{encoding
strategies}: practices that increase the depth and connectivity of
memory traces.

\begin{definition}{Encoding Depth}{encoding-depth}
  The \emph{encoding depth} of an experience $E$ is the number of
  constraint modifications $|\Delta\mathcal{C}|$ it induces — how
  many existing concepts are connected to and modified by the new
  experience. Shallow encoding ($|\Delta\mathcal{C}|$ small) produces
  weak, isolated memories; deep encoding produces strong, well-connected ones.
\end{definition}

Elaborative rehearsal (connecting new information to existing knowledge),
self-explanation (articulating what a concept means in your own terms),
and testing (attempting to retrieve and apply knowledge) all increase
encoding depth. They increase $|\Delta\mathcal{C}|$ by forcing the
new information into contact with existing constraint structures.

\section{Projects}

Projects are extended goal-directed activities that require sustained
access to a specific subset of the knowledge infrastructure. Managing
projects requires managing which parts of the knowledge infrastructure
are currently active — which constraints are in working memory versus
which are in long-term storage.

\begin{definition}{Project as Cognitive Context}{project-context}
  A \emph{project context} is a subset $\mathcal{C}_P \subset \mathcal{C}$
  of the full constraint set, activated specifically for a project $P$.
  The \emph{context switch cost} when moving between projects $P_1$
  and $P_2$ is proportional to $|\mathcal{C}_{P_1} \setminus
  \mathcal{C}_{P_2}|$ — the constraints that must be deactivated and
  reactivated.
\end{definition}

This formalizes the cost of multitasking: each context switch requires
deactivating one project's constraint context and reactivating another's.
Deep work (sustained attention to a single project) minimizes context
switch costs and allows the full constraint context to be active,
enabling the most complex inferences.

\section{Archives}

Archives are the long-term storage tier of the personal knowledge
infrastructure: the repository of constraint modifications that are
important but not currently active. A well-maintained archive has:
\begin{itemize}
  \item High retrieval efficiency (quick access when needed),
  \item Good index structure (easy to find relevant constraints),
  \item Regular maintenance (constraints updated as knowledge evolves),
  \item Appropriate decay management (important constraints retained,
        trivial ones allowed to decay).
\end{itemize}

\begin{namedthm}{The Personal Knowledge Infrastructure Principle}
  An agent's cognitive performance is bounded by the quality of their
  personal knowledge infrastructure: notes (extended memory), encoding
  practices (memory depth), project management (context efficiency),
  and archive maintenance (long-term retention). Optimizing any one
  tier in isolation is less effective than optimizing the interfaces
  between tiers.
\end{namedthm}

\begin{exercise}
  Audit your current personal knowledge infrastructure. For each tier
  (notes, active memory, projects, archive), identify its current
  implementation, its strengths, and its weakest point. Design one
  specific improvement to the interface between two tiers.
\end{exercise}

\begin{exercise}
  Compare two famous intellectual workers known for prolific output:
  Charles Darwin (extensive notebooks, structured correspondence,
  decades-long project management) and Nikola Tesla (claimed to
  work almost entirely in mental simulation with minimal notes).
  What does each approach imply about their personal knowledge
  infrastructure? What are the failure modes of each?
\end{exercise}
